% Part: applied-modal-logics
% Chapter: temporal-logic
% Section: introduction

\documentclass[../../../include/open-logic-section]{subfiles}

\begin{document}

\olfileid{aml}{tl}{int}

\olsection{جان پہچان}

سمے دیاں منطقاں اوہناں گلاں دے بیان ویکھدیاں نیں جیہڑیاں اگے ہون
گیاں یا پہلاں ہو چکیاں نیں۔ سمے دی منطق دا مڈھ بنھن دا سیہرا آرتھر
پرائر نوں دتا جاندا اے۔ اوہ اینوں \emph{فعل دے زمانے دی منطق} آکھدا سی۔
ایتھے سمے دی منطق نوں سمجھاؤن لئی اسیں زیادہ تر پرائر دے اصل امکانی
طریقے اُتے چلاں گے۔ ایس طریقے وچ سمے دے عامل قضیاتی منطق دے بنیادی
ڈھانچے وچ شامل کیتے جاندے نیں، جیہڑا بیان عام طور اُتے فعل دے زمانے
توں بغیر لیندا اے۔

مثلاً قضیاتی منطق وچ میں بیزی ناں دے اک کُتے دی گل کر سکدا آں، جیہڑا
کدی بہندا اے تے کدی نہیں بہندا، جیویں کُتے عام کردے نیں۔ کلاسیکی منطق
وچ ایہ آکھنا تضاد ہووے گا کہ بیزی بیٹھا اے تے بیزی بیٹھا نہیں۔ پر صاف
اے کہ ایہ دونیں وکھ وکھ ویلے سچ ہو سکدے نیں۔ زبان وچ سمے دے عامل شامل
کر کے اسیں ایہ گل سوکھ نال آکھ سکدے آں۔ سمے دے عامل ایہ وی سمجھاؤندے
نیں کہ ایسے استنباط کیوں معتبر نیں جیویں «بیزی نوں انعام وچ کھان والی
کوئی شے یا اک گیند ملے گی» توں «بیزی نوں انعام وچ کھان والی کوئی شے
ملے گی یا بیزی نوں اک گیند ملے گی»۔

سمے دی منطق نال کئی فلسفیانہ سوال اُٹھدے نیں جیہڑے سانوں اک ڈھانچے
دی تھاں دوجا چُنن ول لے جا سکدے نیں۔ مثلاً آؤن والے سمے دا ممکن پر
غیر لازمی بیان اوہ ہوندا اے جیہڑا نہ لازمی اے نہ ناممکن۔ جے اسیں آکھیے
«رچرڈ بھلکے سودا لین دکان جائے گا»، تاں اسیں اجیہی گل آکھ رہے آں
جیہڑی ہجے نہیں ہوئی، تے جیہدی صداقتی قدر اُتے اختلاف ہو سکدا اے۔
ایس گل اُتے وی اختلاف ہو سکدا اے کہ اوہنوں کوئی صداقتی قدر \emph{دتی}
وی جا سکدی اے یا نہیں۔ جے اسیں پکے جبریت وادی ہوئیے، تاں شاید ایہ منّن
وچ کوئی اوکھ نہ ہووے کہ متعلقہ واقعہ ہوون توں پہلاں وی ایہ جملہ سچ اے
یا جھوٹ۔ ہو سکدا اے سانوں اوہدی صداقتی قدر دا ہجے پتا نہ ہووے۔
دوجے پاسے اسیں منّ سکدے آں کہ آؤن والا سماں واقعی کھلا اے، تے اوہدے
ممکن پر غیر لازمی بیاناں دیاں صداقتی قدراں طے نہیں ہوئیاں۔

سمے دی بناوٹ تے اوہدے سُبھاؤ بارے ایہناں وچوں کئی گلّاں ساڈے چنے
نمونیاں تے سمے دی منطق دے ڈھانچیاں وچ پہلے ای شامل ہوندیاں نیں۔
مثلاً اسیں اپنے آپ نوں پچھ سکدے آں کہ ساڈے نمونیاں وچ سماں اک سیدھی
لڑی ہووے، شاخاں کڈھے، یا چکر وچ چلے۔ سانوں ایہ وی چُننا پے سکدا اے
کہ ساڈے سمے دے نمونیاں دا مڈھ تے آخر ہووے یا نہیں، تے سماں وکھ وکھ
پلاں نال وکھایا جائے یا اک لگاتار پھیلاؤ وانگ۔

\end{document}
